\section{Problemstellung}

Eine Nutzung der verschiedenen Cloud-Plattformen bedingt, dass eine grundlegende Struktur existiert auf welche Anwender der Bundesverwaltung aufbauen können. Das BIT verwendet hierfür den etablierten Begriff der Landingzone.

Als Landingzone wird eine durch den Plattformprovider vordefinierten Struktur bezeichnet, in welcher die effektiven Angebote bezogen werden können.

Aus der bisherigen Erfahrung unterschieden sich die Plattformprovider in der Art und Weise wie ihre Plattform konzeptioniert ist. Die Problemstellung ist somit, dass die Plattformen zwar verschieden sind, jedoch im Idealfall homogenisiert werden müssten. Eine Vereinheitlichung ermöglicht es dem BIT einen Leitrahmen zu definieren welcher für die verschiedenen Projekte gilt, welche ihre Vorhaben auf der Plattform laufen lassen.